% 02_why_not_before.tex

\section{Why the Reverse Was Not Immediate}
\label{sec:why-not-before}

The result is not meant as a neglected algebraic sign flip in Spence's model. Since \citet{spence1973job}, the canonical question has been why low types cannot imitate high types. That question fixes the economic object of cost: the signal is costly because it is hard to produce, and it is harder for the low type than for the high type. Once the friction is specified that way, the reverse problem is mostly ruled out by construction. A high type may face a positive cost, but she is not the type for whom the relevant signal is most dangerous.

The first difference here is that the physical nature of the friction changes. The cost is not production cost but reputation exposure. An ambitious signal is a current asset because it can persuade today, but it is also a liability because it leaves a claim available for adverse repricing tomorrow. Under production cost, the high type is protected by lower marginal cost. Under exposure cost, the same high type may carry more reputation-at-risk. The reversal therefore does not come from assuming that high types suddenly dislike success. It comes from changing what the costly part of costly signaling is.

The second reason is dynamic closure. Standard signaling refinements are built to close off-path beliefs for a fixed signaling environment \citep{cho1987signaling,banks1987equilibrium}. In the usual Spence ranking, that closure works in favor of the high type: the high type has the wider profitable-deviation set, so an off-path high signal is assigned to her. The present model adds a state transition that changes the deviation calculus before the off-path belief is closed. The relevant state is not a gradual learning stock. It is a transition from ordinary evaluation to high scrutiny, where yesterday's ambitious action becomes exposed to adverse repricing. Without that state transition, D1 has no reason to reverse its attribution. With it, D1 is not weakened; it is made to operate on the reversed response-set ordering.

The third reason is the limitation of lying-cost interpretations. Models of deception and truth-telling often put a cost directly on lying or a preference directly on honesty \citep{gneezy2005,kartik2009,abeler2019}. Those objects are important, but they treat honesty as a contemporaneous preference or message cost. The mechanism here treats truthfulness and self-accuracy as an asset that can be carried into the future and repriced. A high type's problem is not that she has a larger intrinsic disutility from lying. It is that her exposed self-model has more value to lose if the evaluation state turns hostile.

This distinction also separates the paper from psychological payoff models. Belief-dependent utility and ego concerns already allow beliefs about the self to matter \citep{benabou2002self,koszegi2006ego,battigalli2009dynamic}. The missing step is to connect that self-assessment object to the selection logic of a signaling game. Once self-accuracy becomes a type-indexed reputation asset, the high type can be the party with the smaller profitable-deviation response set. The Spence signal then fails from above: not because low types climb up, but because good types climb down.

These three modeling choices explain why the reverse was not automatic. If costs are production costs, dynamics are only learning or reputation accumulation, and honesty is an exogenous preference, then the standard bottom-up screening question is the natural one. Top-down pooling requires a different closure: a signal that creates future exposure, a state transition that prices that exposure, and a true-type asset that makes the high type more vulnerable to adverse repricing.
